%% ma: L10-B12-0040
%% lop: 10
%% chuong: 5
%% bai: 12
%% dang: 10.12.6
%% loai: TN4
%% muc: TH
%% dap_an: "A"
%% nguon: "chuyen-de-cac-so-dac-trung-cua-mau-so-lieu-khong-ghep-nhom-toan-10-chuong-trinh-moi.pdf (GV. Phan Nhật Linh) (P01), Bài 01 Phần I Câu 21"
%% kien_thuc: "Lớp 10 Bài 12 (số gần đúng, sai số tuyệt đối, độ chính xác, sai số tương đối, quy tròn số)"
%% trang_thai: chinh-thuc
%% ghi_chu: "Đề gốc hỏi dạng chuẩn của số gần đúng (không có trong SGK KNTT, phương án C và D gốc bằng nhau); đã sửa thành hỏi số quy tròn, đáp án không đổi về giá trị (3 214 000). Tự duyệt theo yêu cầu giáo viên 22/10/2026."
\begin{ex}
Số dân của một tỉnh là $a=3\,214\,056$ người với độ chính xác $d=100$ người. Số quy tròn của $a$ là
\choice{\True $3\,214\,000$}{$3\,214\,100$}{$3\,200\,000$}{$3\,210\,000$}
\loigiai{$d=100$ có chữ số khác $0$ đầu tiên ở hàng trăm, nên quy tròn $a$ đến hàng nghìn: $3\,214\,000$. Chọn A.}
\end{ex}
